%%%%%%%%%%%%%%%%%%%%%%%%%%%%%%%%%%%%%% Section 1: Introduction %%%%%%%%%%%%%%%%%%%%%%%%%%%%%%%%%%%%
\section{Introduction} 
\label{sec:Introduction}

As Unmanned Aerial Vehicles (UAVs) operate in increasingly dynamic and complex environments, the demand for reliable navigation \cite{vitale2022autonomous}, including efficient and adaptive path-planning strategies \cite{vitale2024probabilistically}, has grown significantly. Path planning, a critical component of autonomous UAV navigation, determines the optimal path from a starting point to a target region while avoiding obstacles, often optimizing specific mission objectives \cite{lavalle2006planning,papaioannou2023distributed,10925885,papaioannou2025rolling}. This process is central to applications such as emergency response \cite{wan2022accurate, papaioannou20213d, papaioannou2021towards,papaioannou2020coordinated,papaioannou2024synergising}, surveillance \cite{papaioannou2022distributed,wu2020cooperative, papaioannou2019probabilistic, papaioannou2020cooperative,papaioannou2023joint}, and automated inspection \cite{jing2020multi, sun2022automating,papaioannou2022uav,PAPAIOANNOU2024532,papaioannou2023cooperative,papaioannou2023unscented}.

Existing sampling-based path-planning algorithms, such as Rapidly-exploring Random Trees (RRT) \cite{lavalle1998rapidly, kuffner2000rrt}, offer significant advantages, including their ability to handle high-dimensional spaces and their probabilistic completeness, meaning they will eventually find a solution if one exists. However, these methods require careful fine-tuning and often fail to converge to optimal solutions, particularly in complex or cluttered environments. This limitation is critical, as it reduces their reliability, making them less suitable for high-stakes applications like search-and-rescue and disaster response missions, where rapid and dependable solutions are paramount. Recent hybrid approaches combining sampling-based methods with machine learning techniques have shown promise in addressing these shortcomings, offering improved computational efficiency and path quality \cite{arslan2015machine, wang2021deep, kuo2018deep}.

In this direction, this work integrates multimodal large language models (LLMs) with RRT-based path planning to tackle these challenges. The proposed framework, Vision Language Model RRT (VLM-RRT), combines the strengths of sampling-based path planning through RRT with the pattern-matching capabilities and emergent reasoning of LLMs, thereby enhancing autonomous UAV navigation by enabling efficient and robust path planning. Specifically, the proposed approach incorporates a VLM module into the path-planning process to analyze environmental snapshots and dynamically guide the planner to prioritize and sample from specific regions. This reduces redundant exploration and accelerates convergence by biasing the sampling process toward regions more likely to contain feasible and efficient paths, thus improving convergence rates and path quality. The contributions of this work can be summarized as follows:

\begin{itemize}
\item We propose VLM-RRT, a novel framework that augments traditional RRT-based path planning with the advanced reasoning capabilities of Vision Language Models (VLMs). By leveraging a VLM module to analyze environmental snapshots, VLM-RRT enhances navigation decisions by dynamically guiding the sampling process, effectively biasing it toward regions with a higher likelihood of containing optimal paths.
\item Extensive qualitative and quantitative experimental evaluations using OpenAI's GPT-4o and Meta's Llama 3.2 multimodal LLMs demonstrate the effectiveness of the proposed approach in terms of sampling efficiency and path quality compared to the standalone RRT approach.
\end{itemize}

The remainder of this paper is organized as follows. Section \ref{sec:Related_Work} provides background and discusses related work, Section \ref{sec:Problem} formulates the problem, and Section \ref{sec:Approach} details the proposed approach. Finally, Section \ref{sec:Evaluation} evaluates the proposed approach, and Section \ref{sec:Conclusion} concludes the paper.

